\question
Память кольца. Перед дальним патрулём Джон Стюарт освобождает память кольца, сжимая журнал с алфавитом $\{A,B,C,D,E,F\}$. Частоты символов во всём журнале указаны в таблице.
\begin{center}\begin{tabular}{c|rrrrrr}Сигнал&$A$&$B$&$C$&$D$&$E$&$F$\\\hline Частота&15&8&6&5&4&2\end{tabular}\end{center}
\begin{enumerate}
\item Постройте префиксные коды Хаффмана и Шеннона---Фано. При равноценных выборах допустим любой вариант; покажите деревья. Каждым кодом закодируйте фрагмент $\mathrm{CABADEAF}$.
\item Для каждого кода найдите среднюю длину кодового слова и длину кодирования всего журнала. Сравните их с равномерным двоичным кодом минимальной длины для этого алфавита. Служебную таблицу кодов при подсчёте длины не учитывайте.
\end{enumerate}

---------------

Автор -- Олег Эмрих
